% =====================================================================
%  Appendix module: Attribution of the modelled sink to historical
%  carbon, litter inputs and climate
%  Modular subfile -- see main file header.
%
%  STATUS (2026-10-08): DRAFT FOR LORENZO'S REVIEW. Not included in the
%  draft yet. Code: doublechecks/attribution_decomposition.R (numbers),
%  manuscript/figures/build_S15S17_attribution.R (figures S15-S17).
%  Origin: S. Launiainen's total-differential note (Calibration_real_data_
%  transient/documentation/Launiainen_20261008_...pdf). The target was
%  revised on 2026-10-08 (Lorenzo): "historical carbon" must tell the
%  INITIALISATION story, so the sink itself is split, with the initial
%  state as one of the three parts, instead of the year-to-year change in
%  the sink with a stock-feedback term. Open choices are listed at the end.
% =====================================================================
\documentclass[../HIKET_draft_v1.tex]{subfiles}

\begin{document}

\section{Where the modelled sink comes from: the litter rise before 1985, the litter change since, and climate}
\label{app:attribution}

The calibrated models reproduce a soil carbon sink that rises and then slows between
1985 and 2024, and the projection carries the sink forward to 2084. The present
appendix asks how much of that sink, in each year, comes from each of three sources:
the state in which the soil entered the observed period, the litter entering the soil,
and the climate that sets the rate of decomposition. The question is answered for
each of the six models separately, with the same procedure for every model whatever
its number of pools or its form of climate response, and the six answers are then
averaged.

\subsection{The quantity divided and the three parts}

Each model advances the soil carbon stock one year at a time, and the annual sink is
the change in stock over the year,
\begin{equation}
  F_t \;=\; C_t - C_{t-1}.
  \label{eq:attr_sink}
\end{equation}
The sink of every year from 1917 to 2084 is divided into three parts that add up to
it exactly:
\begin{description}
  \item[Litter rise before 1985 (the initialisation)] the part of the sink that the
    soil owes to the state in which it entered 1985. Under the transient
    initialisation (Section~\ref{sec:init}) the soil starts in 1917 in balance with the
    litter of that year, and the litter rises to its 1985 level over the following 68
    years while the climate is held at its 1985--2004 mean. The soil lags behind the
    rising litter, so it enters 1985 below the stock it would hold in balance with the
    litter of 1985, and it keeps gaining carbon afterwards to close that deficit. The
    part is the whole sink of the 1917--1984 period, which by construction comes only
    from the rise of litter, and from 1985 onwards the part of the sink that comes from
    the deficit carried into 1985. The inherited state is therefore itself an input
    effect: the delayed effect of the litter having risen before 1985, separated from
    the next part only by the date of the change.
  \item[Litter change since 1985] the part of the sink that comes from the litter
    input departing, after 1985, from its 1985--1989 level.
  \item[Climate] the part of the sink that comes from the climate of each year
    departing from the 1985--2004 mean. All of a model's climate modifiers are moved
    together, so Yasso15 and Yasso20, which carry three modifiers (for the fast pools,
    the N pool and humus), give a single climate part, comparable with the single
    modifier of the other four models.
\end{description}
The two references, the litter of 1985--1989 and the climate of 1985--2004, are the
ones on which the initialisation itself is built: they are the litter and climate with
which the 1985 soil would be in balance had it no deficit.

\subsection{An exact division}

All six models are linear in carbon: for a given sequence of litter and climate, the
stock in any year is the sum of what remains of the initial stock and what the inputs
have added since. The division uses that property directly, through the models' own
simulations, and needs neither derivatives nor a single-pool approximation.

\paragraph{The litter rise before 1985.} Each model is run from 1985 twice under the same litter
and climate: once from its calibrated initial state and once from the equilibrium state
it would have if the 1917 litter had already been at its 1985 level, which is the
equilibrium start of the counterfactual calibration (Section~\gap{counterfactual
section}). Because the models are linear, the difference between the two sinks is the
free relaxation of the initial deficit, and it is the part due to the litter rise
before 1985:
\begin{equation}
  F^{\mathrm{hist}}_t \;=\; F_t\big(\mathbf{C}_{1985}\big) \;-\; F_t\big(\mathbf{C}^{\mathrm{eq}}_{1985}\big).
\end{equation}
Before 1985 the part is the pre-run sink itself, obtained by replaying the
pre-run year by year with the model's run function.

\paragraph{The litter change since 1985, and climate.} What remains, $F_t(\mathbf{C}^{\mathrm{eq}}_{1985})$,
is the sink of a soil that starts in balance; it is zero in every year if litter and
climate stay at their reference values, and it departs from zero only because they do
not. Writing $F(u, x)$ for that sink under litter $u$ and climate $x$, and the
references as $u_0$ and $x_0$, the two parts are
\begin{align}
  F^{\mathrm{inp}}_t &= \tfrac12\big[F_t(u, x_0) - F_t(u_0, x_0)\big] + \tfrac12\big[F_t(u, x) - F_t(u_0, x)\big], \\
  F^{\mathrm{clim}}_t &= \tfrac12\big[F_t(u_0, x) - F_t(u_0, x_0)\big] + \tfrac12\big[F_t(u, x) - F_t(u, x_0)\big],
\end{align}
with $F_t(u_0, x_0) = 0$. Each part is the effect of changing its own driver, averaged
over the two orders in which the two drivers can be changed (the Shapley division; Shapley,
1953\gap{add Shapley1953 to library.bib if the appendix is kept}). The averaging removes
the need to assign the joint effect of a warmer climate acting on more litter to one
driver or the other: half goes to each. By construction
\begin{equation}
  F^{\mathrm{hist}}_t + F^{\mathrm{inp}}_t + F^{\mathrm{clim}}_t \;=\; F_t
\end{equation}
exactly, in every year and for every model.

Five simulations per plot and parameter set give all three parts: the calibrated start
under actual forcing, the equilibrium start under actual forcing, under reference
climate, under reference litter, and under both references. The last one must stay at
equilibrium and is used as a check.

\subsection{From plots and draws to the national picture}

The division is computed for every plot of the set observed in all three campaigns (310
plots, the basis of every observed rate in the paper), from 1917 through the
historical period 1985--2024 and the projection 2025--2084. The projection is driven
exactly as in the main text: the climate of the last twenty years is repeated and the
litter input is held at its 2024 value. Because the division is additive, national
parts are plain averages of plot parts; the parts of a model are averages over
50 posterior draws (averages rather than medians, because averages preserve the
exact sum); and the six-model summary is the unweighted average of the six models.
Summed over time, the parts give the carbon each source has added to the stock since
1917.

The annual sink follows the weather of each year closely, so the figures are given on
annual values and on a ten-year centred running mean. The running mean is taken
separately before and after 1985, so that no window mixes the pre-run, where litter
and climate parts are zero by construction, with the observed period; at the ends of
each period the window shrinks to the years available. The same weights are applied to
every part, so the smoothed parts still add up exactly to the smoothed sink. The
cumulative view needs no smoothing.

\subsection{Checks}

Three checks are run for every model, plot and parameter set before any number is
kept: the pre-run replayed with the run function reproduces the model's initial state
(largest relative difference $1.4\times10^{-14}$); the equilibrium start stays at
equilibrium under the reference litter and climate (largest drift
$2.3\times10^{-13}$\,t\,C\,ha$^{-1}$\,yr$^{-1}$); and the three parts add up to the sink
(within $2\times10^{-15}$).

\subsection{What the division does and does not say}

The division describes the calibrated models, not the soil: each part is the share of
the modelled sink produced by one source, given the parameters retained by the
calibration. The part due to the litter rise before 1985 is as large as the initial
deficit the calibration chose, through $\sigma_{\mathrm{init}}$, and it decays at the
pace of each model's slow pools, so its size and persistence differ between models for
structural reasons; a model whose $\sigma_{\mathrm{init}}$ is close to one carries
almost no inherited sink. Before 1985 the initialisation holds the climate at its
1985--2004 mean, which is also the climate reference, so the climate part is zero there
by construction: a climate that does not change produces no sink of its own, and its
influence on how fast the soil follows the litter rise is included in the litter part.
Driving the pre-run with observed climate (available from 1961) would change the
initialisation itself and is not attempted here.

The climate part is therefore conditional on the climate of the pre-run, at two
levels. At the level of the division, the climate reference is the climate the model
has assumed since 1917, so the climate part measures the climate to which the
modelled soil has not yet adjusted; another reference, with the initialisation
unchanged, would only add a constant term for the gap between the pre-run climate and
the new reference. The \emph{change} in the climate part between periods is the robust
quantity: its step from $-0.02$ (1985--2004) to $-0.33$\,t\,C\,ha$^{-1}$\,yr$^{-1}$
(2005--2024, six-model average) follows the 1.2\,\textdegree C warming between the two
periods, and moving the reference shifts both values while leaving most of their
difference in place. At the deeper level of the initialisation, the pre-run applies the
1985--2004 climate to the years 1917--1984, which were colder. Under a colder pre-run
climate the soil would have entered 1985 with more carbon, a different deficit and more
warming still to adjust to, so the climate part after 1985 would be larger. The climate
effect attributed here is conditional on that assumption of the transient
initialisation, and testing it requires driving the pre-run with a historical climate
series and recalibrating. The climate part contains both the year-to-year
weather and the long-term warming, measured against the 1985--2004 mean. In the
projection the litter is held at its 2024 level, which is above the 1985--1989
reference, so the litter part stays positive while the soil approaches its new balance.
The litter product includes harvest residues and has no separate natural-mortality
term, so neither can be separated from the litter part.

\subsection{Open choices for review}
\label{app:attribution_open}
\begin{enumerate}
  \item \textbf{References.} The litter of 1985--1989 and the climate of 1985--2004 are
    the references of the initialisation, which keeps the litter-rise part identical
    to the difference between the two starts. Other references (for example the
    1961--1990 climate normal, as in the national inventory) would move sink between
    the climate part and the two litter parts.
  \item \textbf{The pre-run inside the litter-rise part.} The pre-run sink is entirely
    due to the litter rise by construction. Alternatively the period could be shown
    separately, or left out.
  \item \textbf{The cumulative view is held back} (carbon gained since 1917, figure S17).
    Its climate part keeps growing through the projection, which a reader takes for
    future warming; the projection repeats 2005--2024 and contains no further warming,
    so that part is the delayed response to the warming already observed, fading as
    the soils settle. The annual view (S15, S16) carries the message without the risk.
  \item \textbf{Annual or smoothed values} in the figures that enter the paper, and the
    width of the running mean (ten years; twenty would span one full cycle of the
    repeated projection climate).
  \item \textbf{The stock-feedback view} of the original proposal (dividing the
    year-to-year \emph{change} in the sink, with a term for the feedback of the growing
    stock on decomposition) was implemented first and replaced; it can be restored as
    a complement if wanted.
\end{enumerate}

\end{document}
